%% kapa: exam=1 round=28 year=2017 period=1 form=A subject=real_estate subject_ko=부동산학원론 seq=6 number=86
%% source: Q-Net 감정평가사 기출문제 (cst00309) artlSeq=5199603 file=extracted/1008/5199603/2017년 제28회 감정평가사 1차시험 문제지/1교시 A형.pdf
%% conversion: kordoc 4.24.51 -> tools/gen1.py (auto); review=proofread-sonnet; flags=none
\begin{question}{86}
프로젝트 파이낸싱(PF)에 의한 부동산개발에 관한 설명으로 옳지 않은 것은?
\begin{choices}
  \item PF는 부동산개발로 인해 발생하는 현금흐름을 담보로 개발에 필요한 자금을 조달한다.
  \item 일반적으로 PF의 자금관리는 부동산 신탁회사가 에스크로우(Escrow) 계정을 관리하면서 사업비의 공정하고 투명한 자금집행을 담당한다.
  \item 일반적으로 PF의 차입금리는 기업 대출 금리보다 높다.
  \item PF는 위험분담을 위해 여러 이해관계자가 계약관계에 따라 참여하므로, 일반개발사업에 비해 사업진행이 신속하다.
  \item PF의 금융구조는 비소구금융이 원칙이나, 제한적소구금융의 경우도 있다.
\end{choices}
\end{question}
\answer{4}
